\scanpage{1}
% not manuscript: page de titre imprimée de l'Université de Montpellier / IMAG.

\scanpage{2}
% not manuscript: couverture portant « Pages 1-292 + table des matières ».

\scanpage{3}
\margin{I}

La "longue marche" à travers la théorie de Galois...

1. Topos multigaloisiens \hfill 1

2. Appl. aux rev. de topos \hfill 4

3. Variantes \unsure{pro-multigaloisiennes|pro-multigaloisiens} \hfill 6

4. Compléments \unsure{remords}

\figure{[Accolade à droite réunissant les rubriques 3 et 4.]}

\uuline{5.} Introduction du contexte arithmétique
"Conjecture anabélienne" fondamentale \hfill \(7'\)

6. Analyse locale de \((X,S)\) en un \(s\in S\). \hfill \(10'\)

7. Reformulation "bordélique" de la conjecture
(le purgatoire nécessaire...) \hfill 12

\deleted{\unsure{6.}}
8. Réflexion taxonomique (définition du cas, et le purgatoire
\unsure{s'aménage} un peu...) \hfill \(25'\)

\deleted{\unsure{7.}}
9. Structure fonctorielle en les \(s\in S\)
(\unsure{Sections|Sentiers} d'extensions "de deuxième degré" \hfill 34

10. \unsure{Ajustement} des hypothèses, (remords) \hfill 38

11. Conclusions sur les syst. de groupoïdes obtenus à partir de
situations géométriques \hfill 39

12. L'analogue topologique. (Où on se convainc aussi que le bordel
groupoïdal peut s'exprimer complètement, dans le cas considéré,
par les \unsure{groupes ext.} à lacets... \hfill 41

13. Retour au cas arithmétique : formulation "galoisienne" \hfill 53

13\textsuperscript{bis}. Retour sur la notion de groupe à lacets \hfill 56

14. Digression cohomologique (sur le "bordage des trous") \hfill 58

\uuline{15.} Retour sur le cas topologique : orbites critiques
\unsure{au dessus} des scindages d'extensions
Application aux sous-groupes finis de
\unsure{\(\mathrm{Autext}_{\mathrm{lac}}(\pi)\)} (cas discret)
(cf aussi \(\S 18\)) \hfill \(67'\)

16. Bordage et \unsure{forage|fixage} de trous : préliminaires topologiques généraux
\hfill \(79'\)

17. Compléments au \(\S 15\), sur \unsure{op.} de groupes finis
sur des surfaces \hfill 89

\scanpage{4}
\uuline{18.} Forage de trous. \unsure{Applic.} aux s.gr.
finis de \unsure{\(\mathrm{Autext}_{\mathrm{lac}}(\pi)\)}
(cas discret) \hfill 91

\uuline{19.} Tour de Teichmüller \hfill 102

20. Digression : description \unsure{2-homotopique|2-isotopique}
de la cat. des \unsure{2-sphères} topologiques \hfill \(116'\)

21. Les espaces de Teichmüller \hfill 125

\margin{manque \(\S 22\) !}

23. Retour sur les surfaces à groupes (finis) d'opérateurs
("mise en équations" du Pb) \hfill 133

24. Essai de détermination des \illegible{1}. Lien avec les relations
\unsure{\(\Pi_{g,(\nu,\nu+n-1)}^{\Gamma}\)} \(=(1)\),
programme de \deleted{travail} \hfill 142

25. Groupes de Teichmüller "\uline{spéciaux}" \hfill 148

25\textsuperscript{bis}. "Cas des deux groupes d'opérateurs. Retour sur les notations.
\hfill \(153'\)

\uuline{\unsure{26.|25.}} Groupes de Teichmüller profinis. \hfill 161

Discrétification et \unsure{prodiscrétification|prédiscrétification}. Lien avec topos
modulaire de Teichmüller. Conjecture hâtive. \hfill 167

27. Changement de type \((g,\nu)\) : a) bouchage de trous
(et \unsure{diagramme remarquable} p. 174) \hfill 172

28. chgt. de type \((g,\nu)\) : b) passage à un rev. fini
(la conjecture hâtive grince...) \hfill 177

29. Critique de l'approche précédente. (On \unsure{regarde|reprend}
les notions et les conjectures...) \hfill 186

30. Propriétés des \(N_{g,\nu},\Gamma_{g,\nu}\) : a)
\unsure{Propriétés liées} aux sous-groupes finis de Teichmüller.
\hfill 192

31. Digression sur les relèvements d'une action extérieure d'un groupe
fini sur un groupe profini à lacets \hfill 198

\uuline{32.} Retour sur les aspects arithmétiques du bouchage de trous :
relations entre \(\Gamma_{g,\nu}\) et \(\Gamma_{g,\nu-1}\).
(Où on reconstitue, sans le dire \uline{tout le topos étale} d'une
courbe algébrique \inserted{complète} à partir du \(\pi_1\) d'un
ouvert \unsure{auxiliaire}...) \hfill 200

\uuline{33.} Digression topologique : anti-involutions des surfaces
orientées compactes (en laissant de côté d'abord le cas "à trous")
\hfill 221

\uuline{\unsure{35.}} L'injectivité de
\unsure{\(\Gamma_{\mathbb Q}\longrightarrow
\mathrm{Autext}_{\mathrm{lac}}(\widehat{\Pi}_{0,3})\)}
\hfill 249

\unsure{36.} L'isomorphisme
\unsure{\(\Gamma_{\mathbb Q}\simeq\overline{\Pi}_{1,1}\)}
et l'injectivité de
\unsure{\(\Gamma_{\mathbb Q}\longrightarrow
\mathrm{Autext}(\widehat{\Gamma}_{1,1}^{*})
=\mathrm{Autext}(\mathrm{SL}(2,\mathbb Z)^{\wedge})\)}
\hfill 257

\unsure{37.} Modules des courbes elliptiques via Legendre, ou
\unsure{\(M_{1,1}[2]'\simeq U_{0,3}\ne M_{0,4}!\)} \hfill 272

33. Relations entre \(N_{g,\nu},\Gamma_{g,\nu}\) pour \(g\)
variable (étude des rev. finis) \hfill 239

34. Description homotopique profinie de la cat. des courbes alg.
définies sur des \deleted{\unsure{ext.}}
\unsure{sous-ext.|sur-ext.} finies de \(\mathbb Q_0/\mathbb Q\)
\hfill 245

\figure{[À gauche des deux dernières rubriques, accolade prolongée
vers le haut par une flèche dirigée vers la rubrique
« L'injectivité de \(\Gamma_{\mathbb Q}\) ».]}

\scanpage{5}
\margin{\unsure{814|714}\par\unsure{752|252}\par 832}

\scanpage{6}
\mspage{1}
\margin{\unsure{W|VI}}

1. Topos multigaloisiens.

\uline{Prop.} 1.1. Soit \(\mathcal E\) une catégorie.
Conditions équivalentes :

a) \(\mathcal E\) est un topos, et tout objet de \(\mathcal E\)
est \unsure{loc\textsuperscript{t}} constant.

b) \(\mathcal E\) est équivalent à une catégorie \(\widehat C\),
où \(C\) est un groupoïde.

\margin{N.B. On verra plus bas \unsure{qu'on peut choisir}
\unsure{\(I\)|\(J\)} \uline{canoniquement}}

b') \(\exists\) une famille \((G_i)_{i\in I}\) de groupes
et une équivalence de catégories
\[
\mathcal E\simeq\prod\mathrm{Ens}(G_i)
\]

c) Conditions d'exactitude ad-hoc, du type de celles données dans SGA 1...

Dém. \(b\Longleftrightarrow b'\Longrightarrow a\) immédiat.
Pour \(a)\Longrightarrow b)\), je suis moins sûr, peut être faut il
supposer que \(\mathcal E\) est loc connexe, et qu'il a suffisamment
de foncteurs fibres i.e. suffisamment de points.

\uline{Déf. 1.2.} Si les conditions équivalentes b) b') ci-dessus
sont satisfaites, on dit que \unsure{\(C\)|\(\mathcal E\)}
est \unsure{un|une} \deleted{catégorie} \inserted{topos}
multigaloisien (ou une \uline{catégorie} multigaloisienne).

\margin{\uline{Prop. 1.3}
a) Si \(\mathcal E\) est multigaloisien, les topos induits
\(\mathcal E_{/S}\) aussi.
\par
b) Les sommes de topos multigaloisiens (i.e. \deleted{\illegible{1}}
produits de cat. multigaloisiennes) \unsure{itou}.}

\scanpage{7}
\uline{Prop. 1.4} \unsure{Soient|Soit} \(\mathcal E\) un topos,
\(C\) la catégorie des points de \(\mathcal E\) (opposée à la
catégorie des foncteurs fibres sur \(\mathcal E\)). Le foncteur
canonique
\[
\mathcal E\times C^o\longrightarrow(\mathrm{Ens})
\]
induit un foncteur canonique
\[
\mathcal E\longrightarrow\mathrm{Hom}(C^o,(\mathrm{Ens}))
\overset{\mathrm{déf}}{=}\widehat C
\]

Cas \unsure{part.} : (\(\mathcal E\) étant multigaloisien)

a) \(C\) est un groupoïde (appelé "groupoïde fondamental" du topos
mult. gal. \(\mathcal E\), et \unsure{désormais} noté
\(\Pi_1(\mathcal E)\)).

b) \(\mathcal E\longrightarrow\widehat C\) est une équivalence
de catégories.

Un objet \inserted{\(S\)} d'un topos \inserted{\(\mathcal E\)}
est dit 0-connexe s'il est \(\ne\varnothing_{\mathcal E}\)
(il n'est pas objet initial) et s'il est connexe
(i.e. \(S\simeq S'\amalg S''\) implique
\(S'\) ou \(S''\simeq\varnothing_{\mathcal E}\)) --- cela signifie
aussi que le topos induit \(\mathcal E_{/S}\) est 0-connexe
i.e. n'est pas le topos initial ("topos vide", i.e.
\inserted{équivalent à la} catégorie finale) et qu'il est connexe,
i.e. ... On dit que \(S\) est 1-connexe

\scanpage{8}
\mspage{2}

(ou "simplement connexe") s'il est 0-connexe, et si tout
\deleted{\unsure{rev. de \(S\) est}}
\inserted{objet \(S'\) de \(\mathcal E_{/S}\)
\deleted{\illegible{1}}}
loc const. \unsure{est} est \unsure{\deleted{constant}|constant}
--- ce qui ne dépend \inserted{encore} que du topos induit
\(\mathcal E_{/S}\), qui sera dit alors 1-connexe.

\uline{Prop.} 1.5. Soit \(\mathcal E\) un topos multigaloisien,
et soit \(S\) un objet de \(\mathcal E\). Conditions équivalentes

a) \(S\) est 1-connexe

b) \(S\) \unsure{(est)|est} 0-connexe et projectif

c) le foncteur covariant représenté par \(S\)
\[
T\longmapsto\mathrm{Hom}_{\mathcal E}(S,T)
\]
est un foncteur fibre, ou encore
comme il est déjà exact à g.) il commute aux
\(\underset{\longrightarrow}{\lim}\) quelconques

\inserted{N.B.} (il suffit qu'il commute aux sommes et au passage
au quotient...)

d) \(\mathcal E_{/S}\) est équivalent au topos ponctuel

e) (Si \(\mathcal E=\widehat C\), \(C\) un groupoïde) le foncteur
\(S\) sur \(C\) est \inserted{représentable}

\scanpage{9}
\uline{Déf 1.6} On dit alors, parfois, que \(S\) (ou mieux, le topos
\(\mathcal E_{/S}\)) est un \uline{rev.} \uline{universel} du topos
\inserted{multigal.} \(\mathcal E\).

\uline{Prop. 1.7}
\deleted{\unsure{Prenant} \illegible{2}}
\inserted{Soit \(\mathcal E_0\)} la sous-catégorie pleine de
\(\mathcal E\) formée des objets 1-connexes
(de la cat. multigal. \(\mathcal E\)),
\deleted{Alors} \inserted{et \(\Pi_1(\mathcal E)\)}
le groupoïde fondamental \deleted{\unsure{le gr.}}
\inserted{\unsure{p. 1.3}} de \(\mathcal E\).
On a un foncteur canonique
\[
\mathcal E_0\longrightarrow
\unsure{\Pi_1(C)|\Pi_1(\mathcal E)}
\]

(associant à \unsure{\(\forall\)|tt} \(S\in\mathcal E_0\)
le foncteur fibre qu'il représente, ou plutôt le "point" correspondant
de \unsure{\(C\)|\(\mathcal E\)}), qui est
(non seulement pleinement fidèle, mais mieux) une équivalence de
catégories : tout foncteur fibre sur \(\mathcal E\) est représentable
(par un objet 1-connexe) essentiellement unique, comme de juste...).

\scanpage{10}
\mspage{3}
\deleted{\unsure{Déf. 1.8} \illegible{2}}

\uline{Corollaire 1.8.} Soit \(P\) un "point" de \(\mathcal E\)
\inserted{(\unsure{associé} à un foncteur fibre \(F_P\))}.
Il existe un objet 1-connexe \(S\) de \(\mathcal E\) et un relèvement
\(P\xrightarrow{\alpha}\mathcal E_{/S}\)
(i.e. un \(\alpha\in F(S)\)),

\diagram{[Sous \(P\xrightarrow{\alpha}\mathcal E_{/S}\),
\(\mathcal E\) est placé au centre en bas. Deux flèches en pointillé
vont respectivement de \(P\) et de \(\mathcal E_{/S}\) vers
\(\mathcal E\). Entre elles, une flèche vers la droite surmontée
de \(\sim\), \unsure{à double trait|à trait simple}.
Deux crochets latéraux entourent cette partie inférieure du diagramme.]}

et cela détermine \((S,\alpha)\) à isom. unique près.

En fait, \(S\) est l'unique \unsure{objet|\deleted{objet}}
de \(\mathcal E\) qui représente \(F_P\)...

\uline{Déf. 1.9.} On dit que \(S\)
\inserted{(ou \(\mathcal E_{/S}\))} est le
\uline{rev. universel pointé au dessus de \(P\)} déterminé
par le point \(P\).

\uline{Scholie 1.10} Se donner un "point" du \unsure{:}
topos \inserted{multigaloisien} \(\mathcal E\), ou se donner un
rev. universel, revient essentiellement au
\unsure{\(\stackrel{-}{\sim}\)} : chacun détermine l'autre...
